

青青园中葵，朝露待日晞。转眼间我已经在这个学校待了7年，7个春华秋实。时间过得好快，这7年里，我的弟弟从幼儿园的宝宝长到六年级的大孩子，妹妹从小学生变成高考生。也有很多亲密的人离开了我，我的爷爷、姥爷、幺爸、小舅等。

我要感谢我的爸爸妈妈，感谢你们在这么多年给我稳定的支撑、持续的指引、无尽的包容和爱护。无论是什么年纪，你们都在不断学习，不断去面对更大的挑战。我看着你们一步步从筚路蓝缕走到今天，认真地做好每一件小事。是你们给了我不断挑战的决心，不怕从头开始的勇气，以及真诚踏实的处事态度。

感谢亲爱的弟弟，感谢你给我补充“好奇心”的干细胞。很多人长大了之后就会失去对这个世界的好奇心，但是因为有你总是问我“为什么”，让我总是去深入思考一个事情背后的本质原因，你的提问让我养成了这种思考的习惯，让我无论在什么位置，都能像一个孩子一样好奇地去分析这个世界。希望你能健康茁壮成长。

我要感谢我在电子科技大学遇到的所有师长。我觉得是老师塑造了一个学校的风格，让从这里走出去的同学都有了求实求真、大气大为的底色。研究生阶段项目上合作的老师有孙健老师、段景山老师、贾蓉老师、杨宁老师等，他们从技术、表达等方面培养了我严谨认真的做事风格，感谢他们对我的耐心和宽容，感谢他们给我营造的一次次的锻炼机会。感谢吴畏虹老师，他带着我们做了很多项目，以为学生着想的态度帮助每一个学生做学业规划。感谢罗龙老师，她真诚聪明、经验丰富，带着年少的我们初入国际会议的舞台，感谢她的帮助给了我一段难忘的经历。感谢虞红芳老师，感谢她为项目组提供的丰富资源，感谢她为支撑学生自由发展所做的辛勤工作，她以宽容开放的态度在很多次谈话和讲话中激励了我，时间越长越能感觉到虞老师的高瞻远瞩、高屋建瓴与良苦用心。最后我要感谢我的导师孙罡老师，感谢他，感谢他在我毕设工作中的尽心尽力，提出了无微不至的建议，“认真、负责、严谨与追求细节”，是孙老师这次以及之前很多次项目中对我的言传身教；感谢他的课程，进而吸引我加入这个团队，感谢他为团队、为我们所有学生的奔波与辛劳。希望成电蒸蒸日上，团队人才辈出。

感谢我的女朋友赵舒心，这三年因为有你整个学校的色彩都不再一样，有了你之后，我好像就没有了后顾之忧，可以全心全意地向前奔跑。我记得在图书馆，你为我讲解一道道最优化的数学题目，一起准备期末考试；在我难以专注的时候，陪伴我听我梳理源代码；在我疲惫不想睁开眼睛时，我可以安心地让你拉着我走。我们一起经历了很多大大小小的挑战，即使我们的小舟行驶在的是波涛汹涌的大海上，我也没有一丝沮丧，我深深地相信这个小船可以平稳坚定地驶到前方。

感谢这三年里我遇到的每个同学，他们都是非常的优秀。感谢何奕晨和冯伊凡，我们一起在项目和科研上通力合作，忙碌但是充实，又有着悠闲的快乐。感谢伍东旭、郭家豪等同学、冯文佼、谢景昭、马崇喜等师兄师姐、李政廉等师弟师妹的陪伴，希望你们都前途似锦。

感谢我在阿里云实习遇到的每个人，我的leader和mentor们，他们以深厚的技术积累和经验，为我抚平了内心的焦虑、拓展了我未来的路。感谢我足球队的队友们，你们是我的另一个世界。

出生在这个时代从大部分角度来看都是幸运的，感谢有AI的存在，让我这样的学生，能够有比以前高很多倍的效率去学习新的知识，能够探索世界与自己更多的可能性，能够跟随科技更快速地发展。尽管太过强大的AI难免会给人未知的焦虑，我仍然抱有AI乐观主义。

千里之行、始于足下，还需要珍惜生命的每一分钟，珍惜和每个人相处见面的机会，珍惜每一次挑战自己的可能，永远不放弃对这个世界的好奇心。